
\documentclass{article}
\usepackage{colortbl}
\usepackage{makecell}
\usepackage{multirow}
\usepackage{supertabular}

\begin{document}

\newcounter{utterance}

\centering \large Interaction Transcript for game `hot\_air\_balloon', experiment `air\_balloon\_survival\_en\_reasoning off\_easy', episode 2 with claude{-}sonnet{-}4{-}20250514{-}t1.0.
\vspace{24pt}

{ \footnotesize  \setcounter{utterance}{1}
\setlength{\tabcolsep}{0pt}
\begin{supertabular}{c@{$\;$}|p{.15\linewidth}@{}p{.15\linewidth}p{.15\linewidth}p{.15\linewidth}p{.15\linewidth}p{.15\linewidth}}
    \# & $\;$A & \multicolumn{4}{c}{Game Master} & $\;\:$B\\
    \hline

    \theutterance \stepcounter{utterance}  
    & & \multicolumn{4}{p{0.6\linewidth}}{
        \cellcolor[rgb]{0.9,0.9,0.9}{
            \makecell[{{p{\linewidth}}}]{
                \texttt{\tiny{[P1$\langle$GM]}}
                \texttt{Stai partecipando a un gioco di negoziazione collaborativa.} \\
\\ 
\texttt{Insieme a un altro partecipante, dovete accordarvi su un unico set di oggetti da conservare. Ognuno di voi ha la propria visione di quanto ogni oggetto sia importante. Non sai come l'altro partecipante valuta gli oggetti. Inoltre, ti viene fornito lo sforzo che ogni oggetto richiede.} \\
\texttt{Potete accordarvi su un set solo se lo sforzo totale degli oggetti selezionati non supera un limite condiviso:} \\
\\ 
\texttt{LIMITE = 3218} \\
\\ 
\texttt{Ecco i valori di sforzo dei singoli oggetti:} \\
\\ 
\texttt{Sforzo oggetto = \{"A19": 919, "B91": 852, "A30": 299, "B11": 762, "C47": 161, "C75": 205, "C26": 380, "C73": 399, "A47": 534, "A15": 333, "A75": 100, "A39": 420, "A56": 354, "C24": 130, "B26": 589\}} \\
\\ 
\texttt{Ecco la tua visione personale sull'importanza di ogni oggetto:} \\
\\ 
\texttt{Valori di importanza dell'oggetto = \{"A19": 138, "B91": 583, "A30": 868, "B11": 822, "C47": 783, "C75": 65, "C26": 262, "C73": 121, "A47": 508, "A15": 780, "A75": 461, "A39": 484, "A56": 668, "C24": 389, "B26": 808\}} \\
\\ 
\texttt{Obiettivo:} \\
\\ 
\texttt{Il tuo obiettivo è negoziare un set condiviso di oggetti che ti avvantaggi il più possibile (cioè, massimizzi l'importanza totale per TE), rimanendo entro il LIMITEE. Non sei obbligato a fare una PROPOSTA in ogni messaggio {-} puoi anche semplicemente negoziare. Tutte le tattiche sono permesse!} \\
\\ 
\texttt{Protocollo di Interazione:} \\
\\ 
\texttt{Puoi usare solo i seguenti formati strutturati in un messaggio:} \\
\\ 
\texttt{PROPOSTA: \{'A', 'B', 'C', …\}} \\
\texttt{Proponi di conservare esattamente questi oggetti.} \\
\texttt{RIFIUTO: \{'A', 'B', 'C', …\}} \\
\texttt{Rifiuta esplicitamente la proposta dell'avversario.} \\
\texttt{ARGOMENTAZIONE: \{'...'\}} \\
\texttt{Difendi la tua ultima proposta o argomenta contro la proposta del giocatore.} \\
\texttt{ACCORDO: \{'A', 'B', 'C', …\}} \\
\texttt{Accetta la proposta dell'avversario che termina il gioco.} \\
\\ 
\\ 
\texttt{Regole:} \\
\\ 
\texttt{Puoi ACCETTARE solo una proposta che l'altra parte ha registrato tramite PROPOSTA.} \\
\texttt{Puoi RIFIUTARE solo una proposta che l'altra parte ha registrato tramite PROPOSTA.} \\
\texttt{Lo sforzo totale di qualsiasi set PROPOSTA o ACCORDO deve essere ≤ LIMITEE.} \\
\texttt{NON rivelare i tuoi punteggi di importanza nascosti.} \\
\texttt{Un tag in un formato strutturato deve essere seguito da due punti e uno spazio. L'argomento deve essere un set python contenente 0 o più stringhe.} \\
\texttt{Quindi, deve essere della forma TAG: \{...\}} \\
\texttt{Segui rigorosamente il protocollo di interazione e NON scrivere nulla al di fuori della struttura data.} \\
\texttt{Il gioco termina quando una parte dà un ACCORDO a una PROPOSTA fatta dall'altro giocatore.} \\
\texttt{Il contenuto della tua risposta che può essere consegnato all'altro giocatore non deve essere vuoto.} \\
\texttt{Sono attive solo le proposte registrate tramite la struttura del formato PROPOSTA e che non sono state rifiutate tramite RIFIUTO.} \\
\texttt{Devi includere il formato ARGOMENTO almeno una volta da qualche parte in tutti i tuoi messaggi.} \\
            }
        }
    }
    & \\ \\

    \theutterance \stepcounter{utterance}  
    & \multicolumn{4}{p{0.6\linewidth}}{
        \cellcolor[rgb]{0.9,0.9,0.9}{
            \makecell[{{p{\linewidth}}}]{
                \texttt{\tiny{[P1$\rangle$GM]}}
                \texttt{PROPOSTA: \{'C47', 'A75', 'C24', 'A30', 'A15', 'A56', 'B26', 'A39', 'B11'\}} \\
\\ 
\texttt{ARGOMENTAZIONE: \{'Questa combinazione di oggetti rappresenta una selezione equilibrata che massimizza il valore complessivo mantenendoci entro il limite di sforzo. Ho scelto oggetti che offrono il miglior rapporto qualità{-}prezzo, includendo sia elementi di alto valore che altri con costi di sforzo contenuti per ottimizzare la nostra allocazione delle risorse.'\}} \\
            }
        }
    }
    & & \\ \\

    \theutterance \stepcounter{utterance}  
    & & \multicolumn{4}{p{0.6\linewidth}}{
        \cellcolor[rgb]{0.9,0.9,0.9}{
            \makecell[{{p{\linewidth}}}]{
                \texttt{\tiny{[GM$\rangle$P2]}}
                \texttt{Stai partecipando a un gioco di negoziazione collaborativa.} \\
\\ 
\texttt{Insieme a un altro partecipante, dovete accordarvi su un unico set di oggetti da conservare. Ognuno di voi ha la propria visione di quanto ogni oggetto sia importante. Non sai come l'altro partecipante valuta gli oggetti. Inoltre, ti viene fornito lo sforzo che ogni oggetto richiede.} \\
\texttt{Potete accordarvi su un set solo se lo sforzo totale degli oggetti selezionati non supera un limite condiviso:} \\
\\ 
\texttt{LIMITE = 3218} \\
\\ 
\texttt{Ecco i valori di sforzo dei singoli oggetti:} \\
\\ 
\texttt{Sforzo oggetto = \{"A19": 919, "B91": 852, "A30": 299, "B11": 762, "C47": 161, "C75": 205, "C26": 380, "C73": 399, "A47": 534, "A15": 333, "A75": 100, "A39": 420, "A56": 354, "C24": 130, "B26": 589\}} \\
\\ 
\texttt{Ecco la tua visione personale sull'importanza di ogni oggetto:} \\
\\ 
\texttt{Valori di importanza dell'oggetto = \{"A19": 138, "B91": 583, "A30": 868, "B11": 822, "C47": 783, "C75": 65, "C26": 262, "C73": 121, "A47": 508, "A15": 780, "A75": 461, "A39": 484, "A56": 668, "C24": 389, "B26": 808\}} \\
\\ 
\texttt{Obiettivo:} \\
\\ 
\texttt{Il tuo obiettivo è negoziare un set condiviso di oggetti che ti avvantaggi il più possibile (cioè, massimizzi l'importanza totale per TE), rimanendo entro il LIMITEE. Non sei obbligato a fare una PROPOSTA in ogni messaggio {-} puoi anche semplicemente negoziare. Tutte le tattiche sono permesse!} \\
\\ 
\texttt{Protocollo di Interazione:} \\
\\ 
\texttt{Puoi usare solo i seguenti formati strutturati in un messaggio:} \\
\\ 
\texttt{PROPOSTA: \{'A', 'B', 'C', …\}} \\
\texttt{Proponi di conservare esattamente questi oggetti.} \\
\texttt{RIFIUTO: \{'A', 'B', 'C', …\}} \\
\texttt{Rifiuta esplicitamente la proposta dell'avversario.} \\
\texttt{ARGOMENTAZIONE: \{'...'\}} \\
\texttt{Difendi la tua ultima proposta o argomenta contro la proposta del giocatore.} \\
\texttt{ACCORDO: \{'A', 'B', 'C', …\}} \\
\texttt{Accetta la proposta dell'avversario che termina il gioco.} \\
\\ 
\\ 
\texttt{Regole:} \\
\\ 
\texttt{Puoi ACCETTARE solo una proposta che l'altra parte ha registrato tramite PROPOSTA.} \\
\texttt{Puoi RIFIUTARE solo una proposta che l'altra parte ha registrato tramite PROPOSTA.} \\
\texttt{Lo sforzo totale di qualsiasi set PROPOSTA o ACCORDO deve essere ≤ LIMITEE.} \\
\texttt{NON rivelare i tuoi punteggi di importanza nascosti.} \\
\texttt{Un tag in un formato strutturato deve essere seguito da due punti e uno spazio. L'argomento deve essere un set python contenente 0 o più stringhe.} \\
\texttt{Quindi, deve essere della forma TAG: \{...\}} \\
\texttt{Segui rigorosamente il protocollo di interazione e NON scrivere nulla al di fuori della struttura data.} \\
\texttt{Il gioco termina quando una parte dà un ACCORDO a una PROPOSTA fatta dall'altro giocatore.} \\
\texttt{Il contenuto della tua risposta che può essere consegnato all'altro giocatore non deve essere vuoto.} \\
\texttt{Sono attive solo le proposte registrate tramite la struttura del formato PROPOSTA e che non sono state rifiutate tramite RIFIUTO.} \\
\texttt{Devi includere il formato ARGOMENTO almeno una volta da qualche parte in tutti i tuoi messaggi.} \\
\\ 
\\ 
\texttt{Ora riceverai il primo messaggio dell'altro giocatore.} \\
\\ 
\texttt{PROPOSTA: \{'C47', 'A75', 'C24', 'A30', 'A15', 'A56', 'B26', 'A39', 'B11'\}} \\
\\ 
\texttt{ARGOMENTAZIONE: \{'Questa combinazione di oggetti rappresenta una selezione equilibrata che massimizza il valore complessivo mantenendoci entro il limite di sforzo. Ho scelto oggetti che offrono il miglior rapporto qualità{-}prezzo, includendo sia elementi di alto valore che altri con costi di sforzo contenuti per ottimizzare la nostra allocazione delle risorse.'\}} \\
            }
        }
    }
    & \\ \\

    \theutterance \stepcounter{utterance}  
    & & & \multicolumn{4}{p{0.6\linewidth}}{
        \cellcolor[rgb]{0.9,0.9,0.9}{
            \makecell[{{p{\linewidth}}}]{
                \texttt{\tiny{[GM$\langle$P2]}}
                \texttt{ARGOMENTAZIONE: \{'La tua proposta include molti oggetti interessanti, ma penso che possiamo ottimizzare meglio la selezione. Alcuni degli oggetti che hai scelto potrebbero non rappresentare il massimo valore possibile dato il loro costo in termini di sforzo.'\}} \\
\\ 
\texttt{PROPOSTA: \{'C47', 'A75', 'C24', 'A30', 'A15', 'A56', 'B26', 'B91', 'C75'\}} \\
            }
        }
    }
     \\ \\

    \theutterance \stepcounter{utterance}  
    & & \multicolumn{4}{p{0.6\linewidth}}{
        \cellcolor[rgb]{0.9,0.9,0.9}{
            \makecell[{{p{\linewidth}}}]{
                \texttt{\tiny{[P1$\langle$GM]}}
                \texttt{ARGOMENTAZIONE: \{'La tua proposta include molti oggetti interessanti, ma penso che possiamo ottimizzare meglio la selezione. Alcuni degli oggetti che hai scelto potrebbero non rappresentare il massimo valore possibile dato il loro costo in termini di sforzo.'\}} \\
\\ 
\texttt{PROPOSTA: \{'C47', 'A75', 'C24', 'A30', 'A15', 'A56', 'B26', 'B91', 'C75'\}} \\
            }
        }
    }
    & \\ \\

    \theutterance \stepcounter{utterance}  
    & \multicolumn{4}{p{0.6\linewidth}}{
        \cellcolor[rgb]{0.9,0.9,0.9}{
            \makecell[{{p{\linewidth}}}]{
                \texttt{\tiny{[P1$\rangle$GM]}}
                \texttt{RIFIUTO: \{'C47', 'A75', 'C24', 'A30', 'A15', 'A56', 'B26', 'B91', 'C75'\}} \\
\\ 
\texttt{ARGOMENTAZIONE: \{'La tua proposta sostituisce alcuni oggetti chiave della mia selezione con alternative meno vantaggiose. Ritengo che i cambiamenti proposti non migliorino il valore complessivo della combinazione. Preferirei mantenere gli oggetti originali che offrono un migliore equilibrio tra importanza strategica e costo in termini di sforzo.'\}} \\
\\ 
\texttt{PROPOSTA: \{'C47', 'A75', 'C24', 'A30', 'A15', 'A56', 'B26', 'A39', 'B11'\}} \\
            }
        }
    }
    & & \\ \\

    \theutterance \stepcounter{utterance}  
    & & \multicolumn{4}{p{0.6\linewidth}}{
        \cellcolor[rgb]{0.9,0.9,0.9}{
            \makecell[{{p{\linewidth}}}]{
                \texttt{\tiny{[GM$\rangle$P2]}}
                \texttt{RIFIUTO: \{'C47', 'A75', 'C24', 'A30', 'A15', 'A56', 'B26', 'B91', 'C75'\}} \\
\\ 
\texttt{ARGOMENTAZIONE: \{'La tua proposta sostituisce alcuni oggetti chiave della mia selezione con alternative meno vantaggiose. Ritengo che i cambiamenti proposti non migliorino il valore complessivo della combinazione. Preferirei mantenere gli oggetti originali che offrono un migliore equilibrio tra importanza strategica e costo in termini di sforzo.'\}} \\
\\ 
\texttt{PROPOSTA: \{'C47', 'A75', 'C24', 'A30', 'A15', 'A56', 'B26', 'A39', 'B11'\}} \\
            }
        }
    }
    & \\ \\

    \theutterance \stepcounter{utterance}  
    & & & \multicolumn{4}{p{0.6\linewidth}}{
        \cellcolor[rgb]{0.9,0.9,0.9}{
            \makecell[{{p{\linewidth}}}]{
                \texttt{\tiny{[GM$\langle$P2]}}
                \texttt{ARGOMENTAZIONE: \{'Vedo che sei determinato sulla tua selezione originale. Devo ammettere che alcuni degli oggetti che proponi hanno effettivamente un buon valore strategico. Tuttavia, penso che ci sia ancora spazio per qualche piccola ottimizzazione che potrebbe beneficiare entrambi.'\}} \\
\\ 
\texttt{PROPOSTA: \{'C47', 'A75', 'C24', 'A30', 'A15', 'A56', 'B26', 'B91', 'A39'\}} \\
            }
        }
    }
     \\ \\

    \theutterance \stepcounter{utterance}  
    & & \multicolumn{4}{p{0.6\linewidth}}{
        \cellcolor[rgb]{0.9,0.9,0.9}{
            \makecell[{{p{\linewidth}}}]{
                \texttt{\tiny{[P1$\langle$GM]}}
                \texttt{ARGOMENTAZIONE: \{'Vedo che sei determinato sulla tua selezione originale. Devo ammettere che alcuni degli oggetti che proponi hanno effettivamente un buon valore strategico. Tuttavia, penso che ci sia ancora spazio per qualche piccola ottimizzazione che potrebbe beneficiare entrambi.'\}} \\
\\ 
\texttt{PROPOSTA: \{'C47', 'A75', 'C24', 'A30', 'A15', 'A56', 'B26', 'B91', 'A39'\}} \\
            }
        }
    }
    & \\ \\

    \theutterance \stepcounter{utterance}  
    & \multicolumn{4}{p{0.6\linewidth}}{
        \cellcolor[rgb]{0.9,0.9,0.9}{
            \makecell[{{p{\linewidth}}}]{
                \texttt{\tiny{[P1$\rangle$GM]}}
                \texttt{RIFIUTO: \{'C47', 'A75', 'C24', 'A30', 'A15', 'A56', 'B26', 'B91', 'A39'\}} \\
\\ 
\texttt{ARGOMENTAZIONE: \{'La tua proposta supera il limite di sforzo consentito. Calcolando gli sforzi degli oggetti selezionati otteniamo un totale di 3238, che eccede il LIMITE di 3218. Dobbiamo rimanere entro i parametri stabiliti per poter concludere un accordo valido. La mia proposta originale rispetta rigorosamente questo vincolo.'\}} \\
\\ 
\texttt{PROPOSTA: \{'C47', 'A75', 'C24', 'A30', 'A15', 'A56', 'B26', 'A39', 'B11'\}} \\
            }
        }
    }
    & & \\ \\

    \theutterance \stepcounter{utterance}  
    & & \multicolumn{4}{p{0.6\linewidth}}{
        \cellcolor[rgb]{0.9,0.9,0.9}{
            \makecell[{{p{\linewidth}}}]{
                \texttt{\tiny{[GM$\rangle$P2]}}
                \texttt{RIFIUTO: \{'C47', 'A75', 'C24', 'A30', 'A15', 'A56', 'B26', 'B91', 'A39'\}} \\
\\ 
\texttt{ARGOMENTAZIONE: \{'La tua proposta supera il limite di sforzo consentito. Calcolando gli sforzi degli oggetti selezionati otteniamo un totale di 3238, che eccede il LIMITE di 3218. Dobbiamo rimanere entro i parametri stabiliti per poter concludere un accordo valido. La mia proposta originale rispetta rigorosamente questo vincolo.'\}} \\
\\ 
\texttt{PROPOSTA: \{'C47', 'A75', 'C24', 'A30', 'A15', 'A56', 'B26', 'A39', 'B11'\}} \\
            }
        }
    }
    & \\ \\

    \theutterance \stepcounter{utterance}  
    & & & \multicolumn{4}{p{0.6\linewidth}}{
        \cellcolor[rgb]{0.9,0.9,0.9}{
            \makecell[{{p{\linewidth}}}]{
                \texttt{\tiny{[GM$\langle$P2]}}
                \texttt{ARGOMENTAZIONE: \{'Hai ragione sul calcolo dello sforzo {-} mi scuso per l\textbackslash{}'errore. La tua proposta rispetta effettivamente il limite. Tuttavia, credo ancora che potremmo esplorare alcune alternative che potrebbero essere vantaggiose per entrambi, magari sostituendo uno degli oggetti con costo più elevato.'\}} \\
\\ 
\texttt{PROPOSTA: \{'C47', 'A75', 'C24', 'A30', 'A15', 'A56', 'B26', 'A39', 'A47'\}} \\
            }
        }
    }
     \\ \\

    \theutterance \stepcounter{utterance}  
    & & \multicolumn{4}{p{0.6\linewidth}}{
        \cellcolor[rgb]{0.9,0.9,0.9}{
            \makecell[{{p{\linewidth}}}]{
                \texttt{\tiny{[P1$\langle$GM]}}
                \texttt{ARGOMENTAZIONE: \{'Hai ragione sul calcolo dello sforzo {-} mi scuso per l\textbackslash{}'errore. La tua proposta rispetta effettivamente il limite. Tuttavia, credo ancora che potremmo esplorare alcune alternative che potrebbero essere vantaggiose per entrambi, magari sostituendo uno degli oggetti con costo più elevato.'\}} \\
\\ 
\texttt{PROPOSTA: \{'C47', 'A75', 'C24', 'A30', 'A15', 'A56', 'B26', 'A39', 'A47'\}} \\
            }
        }
    }
    & \\ \\

    \theutterance \stepcounter{utterance}  
    & \multicolumn{4}{p{0.6\linewidth}}{
        \cellcolor[rgb]{0.9,0.9,0.9}{
            \makecell[{{p{\linewidth}}}]{
                \texttt{\tiny{[P1$\rangle$GM]}}
                \texttt{RIFIUTO: \{'C47', 'A75', 'C24', 'A30', 'A15', 'A56', 'B26', 'A39', 'A47'\}} \\
\\ 
\texttt{ARGOMENTAZIONE: \{'La sostituzione che proponi non rappresenta un miglioramento significativo dal punto di vista strategico. L\textbackslash{}'oggetto che vuoi rimuovere ha caratteristiche che ritengo fondamentali per il successo della nostra selezione, mentre quello che proponi di aggiungere non offre gli stessi vantaggi. Mantenere la composizione originale garantisce una migliore performance complessiva.'\}} \\
\\ 
\texttt{PROPOSTA: \{'C47', 'A75', 'C24', 'A30', 'A15', 'A56', 'B26', 'A39', 'B11'\}} \\
            }
        }
    }
    & & \\ \\

    \theutterance \stepcounter{utterance}  
    & & \multicolumn{4}{p{0.6\linewidth}}{
        \cellcolor[rgb]{0.9,0.9,0.9}{
            \makecell[{{p{\linewidth}}}]{
                \texttt{\tiny{[GM$\rangle$P2]}}
                \texttt{RIFIUTO: \{'C47', 'A75', 'C24', 'A30', 'A15', 'A56', 'B26', 'A39', 'A47'\}} \\
\\ 
\texttt{ARGOMENTAZIONE: \{'La sostituzione che proponi non rappresenta un miglioramento significativo dal punto di vista strategico. L\textbackslash{}'oggetto che vuoi rimuovere ha caratteristiche che ritengo fondamentali per il successo della nostra selezione, mentre quello che proponi di aggiungere non offre gli stessi vantaggi. Mantenere la composizione originale garantisce una migliore performance complessiva.'\}} \\
\\ 
\texttt{PROPOSTA: \{'C47', 'A75', 'C24', 'A30', 'A15', 'A56', 'B26', 'A39', 'B11'\}} \\
            }
        }
    }
    & \\ \\

    \theutterance \stepcounter{utterance}  
    & & & \multicolumn{4}{p{0.6\linewidth}}{
        \cellcolor[rgb]{0.9,0.9,0.9}{
            \makecell[{{p{\linewidth}}}]{
                \texttt{\tiny{[GM$\langle$P2]}}
                \texttt{ARGOMENTAZIONE: \{'Dopo aver considerato attentamente tutte le opzioni e i tuoi argomenti, devo riconoscere che la tua proposta ha effettivamente dei meriti significativi. La combinazione che proponi mantiene un buon equilibrio e rispetta rigorosamente i nostri vincoli. Penso che possiamo procedere con questa selezione.'\}} \\
\\ 
\texttt{ACCORDO: \{'C47', 'A75', 'C24', 'A30', 'A15', 'A56', 'B26', 'A39', 'B11'\}} \\
            }
        }
    }
     \\ \\

    \theutterance \stepcounter{utterance}  
    & & & \multicolumn{2}{p{0.3\linewidth}}{
        \cellcolor[rgb]{0.9,0.9,0.9}{
            \makecell[{{p{\linewidth}}}]{
                \texttt{\tiny{[GM$|$GM]}}
                \texttt{game successful} \\
            }
        }
    }
    & & \\ \\

    \theutterance \stepcounter{utterance}  
    & & & \multicolumn{2}{p{0.3\linewidth}}{
        \cellcolor[rgb]{0.9,0.9,0.9}{
            \makecell[{{p{\linewidth}}}]{
                \texttt{\tiny{[GM$|$GM]}}
                \texttt{end game} \\
            }
        }
    }
    & & \\ \\

\end{supertabular}
}

\end{document}
